\section{Methods}
\label{sec:methods}

\subsection{Study area and historical cores}
\label{sec:extent}

The study area was defined as the southern African cheetah range \parencite{weise2017distribution} plus a buffered analysis extent. All layers were aligned to a fixed equal-area projection and a \cellsize\ snap grid, so every raster was directly comparable cell for cell. Cores were derived from the historical confirmed-range baseline and held constant across 2012, 2016, 2020, and 2024. The final network consisted of 23 cores and 45 selected neighboring core pairs.

% Working resolution belongs with the study area, not in a separate subsection.
\label{sec:scale}

I used a 1-km working resolution because it represented a compromise between the rougher 10-km occurrence baseline and the finer environmental covariates, which ranged from 30 m to 1 km. This resolution was fine enough to identify meaningful points in the landscape while avoiding the implication of a level of accuracy that the occurrence evidence alone could not provide.

% Core definition belongs in this same subsection.
\label{sec:cores}

I defined a core area as a contiguous patch of cells that were confirmed presence, with a density of $\geq$ 0.51 cheetahs per 100 $km^2$. The final network consisted of 23 cores, which were held constant across all four temporal snapshots to avoid inventing core movement in the absence of comparable range reassessment.

I represented these 23 core areas as fixed nodes across all four temporal snapshots (2012, 2016, 2020 and 2024). I took each core area and calculated the Euclidean distance to the nearest neighbor. This yielded 45 straight-line links. These straight-line links defined the lines used to get the least cost paths based on the added resistance surface rasters.

\subsection{Landscape data and resistance surfaces}
\label{sec:covariates}

The predicted temporal resistance values came from the continuous rasters (built-up fraction (human development) from GHSL, nighttime lights from VIIRS, and VCF (vegetation) from MODIS). Categorical land cover was used for masking and description only, because annual categorical features can switch classes between years for reasons unrelated to overall landscape change.

I also used the static covariates (roads, livestock, terrain, hydrography, fences, protected areas, and core definitions). The static covariates did not contribute \textit{temporal} signal to the resistance surface; however, they did contribute \textit{spatial} signal to the resistance surface. The static covariates were used to describe the landscape and to provide context for the dynamic covariates.

% Explain how the input layers become resistance here.
\label{sec:resistance}

Resistance values were assigned to represent the relative difficulty of movement through each landscape feature. The terrain component was derived from slope, with steeper slopes assigned higher resistance values. I compared alternative vegetation weights because vegetation cover can significantly influence cheetah movement, \TODO{add a cite for this} and different weighting schemes allowed for sensitivity analysis of just this factor. the final balanced map \TODO{figure?} had a weight of 60\% for human pressure, 20\% for vegetation, and 20\% for terrain. Within human pressure, built-up land, roads, and livestock each received 30\%, while nighttime lights received 10\%. Available fence data was incorporated into the resistance model, with fence cells increasing the combined resistance by a multiplier from 1 to 25.

\begin{outline}

Terrain component. The DEM was not used directly as resistance; slope was derived from
it. Mean slope carries the primary terrain surface, maximum slope retained for
sensitivity and context. Terrain resistance function: slope $\le 10^\circ$ low
resistance; $10$ to $30^\circ$ linearly increasing; $\ge 30^\circ$ capped high;
resistance rescaled to 1 to 10.

Original primary temporal surface. Anthropogenic component 80\% of total, terrain 20\%.
Within the anthropogenic component: built-up 30\%, roads 30\%, livestock 30\%,
nighttime lights 10\%.

Three internal weighting models compared: primary, equal, infrastructure-emphasis.

Vegetation correction. The first temporal model omitted VCF from the resistance
calculation, which is a mismatch with the predeclared protocol. Three further
sensitivity scenarios were added, keeping the original anthropogenic-only surfaces
intact rather than replacing them:
\begin{itemize}
  \item vegetation-low: 10\% vegetation, 70\% anthropogenic, 20\% terrain
  \item vegetation-balanced: 20\% vegetation, 60\% anthropogenic, 20\% terrain
  \item vegetation-high: 40\% vegetation, 40\% anthropogenic, 20\% terrain
\end{itemize}
The vegetation proxy uses continuous VCF fields to represent sparse and bare cover. It
does not claim to measure fine-scale habitat edge or observed movement preference, and
should be described that way.

Reconciliation still needed: these six weighting models are not the same object as the
predeclared RES-01 to RES-06 scenario set. Either map each one onto a predeclared
scenario ID or state plainly in the text that the implemented set differs and why.
\end{outline}

Resistance values were assigned to represent the relative difficulty of movement through each landscape feature. The terrain component was derived from slope, with steeper slopes assigned higher resistance values. This had to be done carefully however, given cheetahs do have an affinity for some sloped terrains \TODO{PLAIN-LANGUAGE COMPARISON: \textcite{moqanaki2016iran} treated slopes near 15 percent as easiest for movement and assigned more resistance above or below that point. Explain that this study instead used chosen slope ranges as a sensitivity assumption because no local GPS data were available to estimate the exact response.}. I compared alternative vegetation weights because vegetation cover can significantly influence cheetah movement, and different weighting schemes allowed for sensitivity analysis of this factor.

\subsection{Connectivity modelling and temporal comparison}
\label{sec:connectivity}
% Sentence starters: "For each selected core pair, I calculated ..." / "The straight-line link was used to define ..." / "The least-cost path could depart from that line when ..."
\TODO{PLAIN-LANGUAGE WRITING GUIDE.
Least-cost paths: for each starting core, ArcGIS calculated the accumulated cost of
moving through every reachable cell. A path was then drawn from each connected core
back to the starting core by following the lowest accumulated cost. Reusing one cost
calculation for several destinations gives the same pairwise paths with fewer repeated
calculations.
Circuitscape: Circuitscape 5.17.1 \parencite{anantharaman2020circuitscape} treated the
landscape like an electrical network and spread current across all possible routes
between each pair of cores. All 253 core pairs were solved in each of four years, for
1,012 pair comparisons, with no failures.
Why cores were removed from reported current: cores are where modeled current enters
and leaves, so they automatically receive the highest values. Removing the core cells
after calculation lets the maps focus on land between cores. This improves what is
reported, but the cores still influenced the calculation.
Why a core-to-core method was used: the study began with defined core areas, so the
question was how those known areas connect. An all-direction method such as Omniscape
\parencite{landau2021omniscape} answers a different question.}


% Current classes are part of reporting the connectivity model, not a separate corridor definition.
\TODO{PLAIN-LANGUAGE WRITING GUIDE: do not draw a hard boundary and call everything
inside it a confirmed corridor. Report the full continuous current surface. Percentile
classes are used only to make the maps easier to read and to compare protected-area
coverage. Percentiles are calculated after core cells are removed, and the class
boundaries were chosen before looking at the protected-area result.}

% Network importance belongs alongside the connectivity calculations.
\label{sec:graph}
\TODO{PLAIN-LANGUAGE WRITING GUIDE.
What was measured: edge betweenness counts how often a link lies on the shortest
network route between pairs of cores. A high value means the link helps hold many parts
of the modeled network together. The value was calculated on the 45-link network and
divided by the 253 possible core pairs.
How links were chosen: the network combined each core's three nearest neighbors with
a minimum spanning tree, which is the smallest set of links needed to keep the core
network connected.
Why rankings are uncertain: allowing two, three, four, or five nearest neighbors
created 31, 45, 61, or 77 links, and the highest-ranked links changed. Link 17--24 was
the only bridge under the three-neighbor choice; removing a bridge separates the
network into two pieces.
What was dropped: PC and dPC \parencite{saura2007pc} require an assumed relationship
between distance and dispersal. No local dispersal telemetry was available, so those
more detailed-looking numbers were not calculated. Cite \textcite{keeley2021metrics}
when explaining that the metric should match the research question.}

% The guide above records why PC/dPC were not calculated; no dispersal placeholder is needed.
\label{sec:temporal}
% Sentence starters: "I compared the same core pairs across ..." / "Absolute change was calculated as ..." / "Proportional change was calculated as ..."
\TODO{PLAIN-LANGUAGE WRITING GUIDE: compare the same core pairs, model assumptions,
and 1-km grid in 2012, 2016, 2020, and 2024. This keeps the comparison consistent so
the modeled differences come from the input layers that changed through time. Define
absolute change as later value minus earlier value. Define percent change as that
difference divided by the earlier value, multiplied by 100.}


\subsection{Protected-area coverage}
\label{sec:pa}
\TODO{PLAIN-LANGUAGE WRITING GUIDE.
Question: what share of the modeled connecting landscape outside the cores falls inside
protected areas?
Data preparation: use the August 2026 WDPA release
\parencite{unepwcmc2026wdpa}. Keep terrestrial and coastal protected-area polygons,
remove proposed areas, points, and marine-only areas, and merge overlaps. This left
2,301 records in the study extent.
Comparison area: 31.26 percent of the usable outside-core area was protected. Cells
with zero current were excluded because they could never enter a high-current group.
Why no p-value: the answer changed from significant to not significant when the
assumed spatial block size changed. Report the descriptive percentages and note that
the same general pattern appeared in all four years and under several core-buffer
sizes.}

\subsection{Sensitivity checks and link screening}
\label{sec:uncertainty}
% Sentence starters: "I considered a result robust when ..." / "A link was classified as sensitive when ..." / "Locations with low agreement were treated as ..."
\TODO{PLAIN-LANGUAGE WRITING GUIDE.
What "robust" means here: a link had to pass three checks. First, both alternative
weightings had to keep at least 80 percent of the route within 5 km of the main route;
34 of 45 passed. Second, the fence scenarios could not substantially change the link;
this removed two more from the group. Third, the direction of change had to stay the
same under all vegetation weightings. The two links that failed this last check had
already been removed. The final groups were 32 main-priority links and 13 uncertain
links.
Important distinction: the 80-percent-within-5-km rule is not the cell-by-cell agreement
test proposed in the original protocol, even though both use the number 80 percent.
What was not completed: alternative core definitions, 2-km and 5-km grids, and a
cell-by-cell agreement map.
Separate location uncertainty from cost uncertainty: 49 of 180 route-and-year
comparisons had less than 80 percent overlap within 1 km, including 21 main-priority
links. Link 17--24 had no overlap in 2020. This does not mean the connection was
unimportant; it means the exact mapped route was uncertain.}

\begin{outline}
\begin{itemize}
  \item 90 weight comparisons evaluated.
  \item 34 of 45 core pairs robust under the tested alternative weights, 11 sensitive.
  \item Fence scenarios changed 5 links: 1--13, 1--17, 9--13, 13--17, 27--30.
  \item 32 links classified primary-priority eligible; 13 retained as uncertainty and
        data-collection priorities.
\end{itemize}
Two things to resolve before writing: state the numeric robustness rule actually applied
and whether it matches the predeclared $k/n = 0.8$; and reconcile the 34 robust pairs
with the 32 primary-priority-eligible links, since those are different counts and a
reader will notice.
\end{outline}

% Limited occurrence comparison: a check within the sensitivity section, not independent validation.
\label{sec:validation}
\TODO{PLAIN-LANGUAGE WRITING GUIDE.
What was planned: compare the model with separate Limpopo tracking data and Namibia
camera-trap data.
What was actually done: compare early and late occurrence records from the same
archive used to create the cores. Because the data are not independent, call this a
limited plausibility check rather than validation.
Numbers available: after thinning to one record per 10-km cell, 2010--2012 had 413
cells and 2013--2016 had 450. Their overlap score was 0.239. There are no records after
2016, so this check says nothing about 2020 or 2024.
Decision still needed: either complete the planned outside-data check or clearly state
that it was not completed.}

% Explain the link-level screening output here, without repeating the robustness counts.
\label{sec:mcda}
\TODO{PLAIN-LANGUAGE WRITING GUIDE.
What was planned: combine six measurements into one score for every 10-km cell and
test three sets of weights.
What was built instead: a link-level table that keeps each kind of evidence separate.
Its columns describe how important the link is to the network, how its modeled cost
changed, how stable its mapped route was, whether fences affected it, and how much of
it is protected.
Why this is more honest: combining everything into one score would require deciding,
for example, how much route instability is equal to a change in protection. No
evidence-based exchange rate was available. State clearly that no combined monitoring
score was calculated.}

\subsection{Software and reproducibility}
\label{sec:software}
\TODO{PLAIN-LANGUAGE WRITING GUIDE: list the software without explaining every
technical setting in the main text. Verified versions are ArcGIS Pro
\arcgisversion\ (build \arcgisbuild), ArcGIS Python \pythonversion, Julia
\juliaversion, and Circuitscape.jl \circuitscapeversion; see Supplementary
Table~\ref{tab:software}. Google Earth Engine supplied and aligned the satellite data.
Most ArcGIS processing used Python scripts rather than ModelBuilder. Add the final
archive link and state which source datasets cannot legally be redistributed.}

\begin{outline}
ArcGIS Pro Spatial Analyst via arcpy
(\texttt{CostDistance}, \texttt{CostPathAsPolyline}); Google Earth Engine for SRTM,
GHSL, VIIRS, MODIS MCD12Q1 and VCF acquisition and harmonisation. Record versions.
\end{outline}
